\documentclass[12pt,a4paper]{article}
\usepackage[utf8]{inputenc}
\usepackage[T1]{fontenc}
\usepackage[portuguese]{babel}
\usepackage{geometry}
\usepackage{enumitem}
\usepackage{booktabs}
\usepackage{array}
\usepackage{longtable}
\usepackage{lmodern}
\geometry{margin=2.4cm}
\setlist[itemize]{leftmargin=1.5em, itemsep=0.2em, topsep=0.2em}
\setlist[enumerate]{leftmargin=1.8em}
\setlength{\parindent}{0pt}
\setlength{\parskip}{0.5em}

% Persona: nome, idade e papel, contexto, objetivos, frustrações, ligação ao produto
\newcommand{\persona}[6]{%
  \subsection{#1}
  \textbf{#2}\\
  #3

  \textbf{O que quer.} #4

  \textbf{O que a incomoda hoje.} #5

  \textbf{No sistema.} #6
}

% User story: id, área, persona, quero, para, prioridade proposta, critérios
\newcommand{\us}[7]{%
  \paragraph{#1 \textnormal{(#2, #6)}}
  Como \textit{#3}, quero \textit{#4}, para \textit{#5}.
  \begin{itemize}
  #7
  \end{itemize}
}

\title{Engenharia de Software\\Sprint 1 --- Análise do Produto}
\author{}
\date{17/10/2026 \quad|\quad Sprints 1 e 2 \quad|\quad Feature Request \& Product Roadmap Management System}
\begin{document}
\maketitle

\section*{Estado deste documento}

Este é o documento de base do Sprint 1 e vai ser revisto pela equipa. Contém a visão do produto, os papéis que o enunciado nomeia, as sete áreas de funcionalidade, as personas, os cenários, as user stories com critérios de aceitação, os requisitos não funcionais e as interfaces candidatas.

Duas coisas têm de ficar claras:

\begin{itemize}
\item As personas, os cenários, as histórias e as prioridades são uma \textbf{primeira proposta}. O enunciado não traz nenhuma persona escrita. Partiram do texto do enunciado (\texttt{SE\_Project.pdf}, secção \textit{Project Proposal}) e foram escritas com apoio de IA (Claude, via Claude Code). As pessoas são inventadas, só os papéis e as funcionalidades vêm do enunciado.
\item O enunciado obriga a levantar os requisitos com o Character.ai. Esse levantamento \textbf{ainda não foi feito}. Quando for, os resultados entram aqui para confirmar, corrigir ou cortar o que está escrito. Até lá, nada disto está validado.
\end{itemize}

A prioridade de cada história é uma proposta (método MoSCoW). Quem decide a ordem do backlog é o Product Owner, depois de ouvir a equipa.

\section{Visão e âmbito}

\textbf{Nome.} Development of a Feature Request \& Product Roadmap Management System. Em português: sistema para gerir pedidos de funcionalidades e o roadmap de um produto.

\textbf{Problema.} Um produto de software muda com o feedback dos clientes, com o que os stakeholders pedem, com o mercado e com a tecnologia. Muita gente ainda junta esses pedidos em folhas de cálculo, emails, trackers genéricos ou chats. O resultado é pedidos repetidos, prioridades que não batem certo, pouca ligação entre o que se pediu e o que se fez, e pouca visibilidade de como o produto evolui.

\textbf{O que vamos construir.} Um sistema, com back-end e front-end, em que gestores de produto, clientes, programadores e gestores de projeto trabalham no mesmo sítio, desde o pedido de uma funcionalidade até à release.

\textbf{O que o enunciado exige à volta disto.}
\begin{itemize}
\item Back-end e front-end (o enunciado recomenda Django, React ou Node.js).
\item Pelo menos 8 interfaces para além do login, do registo e do CRUD básico.
\item Arquitetura justificada, com wireframes e mockups.
\item Requisitos levantados com o Character.ai.
\item Pelo menos uma ferramenta de IA para programar.
\item Jira para organizar tudo, com acesso dos docentes.
\item Um Product Owner e um Team Captain (Scrum Master) por sprint, com rotação.
\end{itemize}

\textbf{Fora do âmbito, por agora.} Integração com ferramentas externas (por exemplo importar pedidos do email ou de um tracker) é permitida pelo enunciado, mas não é pedida. Fica para depois do núcleo estar feito.

\section{Quem o enunciado nomeia}

Não são personas. São os papéis que o texto nomeia. As personas da secção seguinte dão rosto a estes papéis.

\begin{itemize}
\item \textbf{Gestores de produto} (product managers). Acompanham a evolução do produto.
\item \textbf{Clientes} (customers). A origem de parte dos pedidos.
\item \textbf{Programadores} (software developers). Recebem trabalho e implementam.
\item \textbf{Gestores de projeto} (project managers). Acompanham o andamento.
\item \textbf{Product Owners.} O enunciado diz que a priorização manual é feita por eles.
\item \textbf{Utilizadores que votam.} O enunciado pede um mecanismo de votação e não diz quem vota. Assumimos que são os clientes e os utilizadores do produto cujos pedidos o sistema gere (a confirmar).
\end{itemize}

\section{Funcionalidades}

São as sete áreas do enunciado, na ordem em que o enunciado as escreve. Essa ordem não é uma prioridade.

\begin{description}
\item[Feature Management.] Submissão de pedidos; categorização e etiquetas; edição e versões; pesquisa e filtros.
\item[Prioritization.] Priorização manual pelos Product Owners; votação pelos utilizadores; sugestões assistidas por IA; estimativa de impacto e de esforço.
\item[Product Backlog Management.] Criação do product backlog; definição do sprint backlog; estado das funcionalidades; atribuição a equipas de desenvolvimento.
\item[Product Roadmap.] Roadmap visual; planeamento de releases; milestones; dependências entre funcionalidades.
\item[Requirements Engineering.] Geração automática de user stories a partir dos pedidos; rastreabilidade entre funcionalidades, user stories e implementação; gestão de critérios de aceitação.
\item[Collaboration.] Comentários e discussões; notificações; menções; espaços de trabalho por equipa.
\item[Dashboards and Analytics.] Pedidos submetidos; aprovados e rejeitados; funcionalidades por release; lead time; cycle time; carga das equipas; progresso do roadmap.
\end{description}

\section{Personas}

Cinco proto-personas, uma por papel principal. Os nomes e as histórias são inventados para ajudar a decidir. O formato segue o das aulas: quem é, o que quer, o que a incomoda.

\persona{Marta Lopes}{34 anos, product manager numa empresa de software B2B}
{Gere um produto usado por cerca de 40 empresas clientes. Recebe pedidos por email, por chat e em reuniões, e guarda-os numa folha de cálculo que só ela percebe. Tem de apresentar o plano do produto à direção de três em três meses.}
{Ver todos os pedidos num só sítio, perceber quais se repetem, decidir com dados o que entra em cada release e mostrar um roadmap que a direção e os clientes entendam.}
{Pedidos duplicados, prioridades decididas por quem insiste mais, e não conseguir explicar porque é que uma funcionalidade ficou de fora.}
{Utilizadora principal do roadmap, das releases, das dependências e do dashboard. Revê a priorização.}

\persona{Ricardo Almeida}{41 anos, Product Owner de uma equipa Scrum}
{Trabalha com uma equipa de seis programadores em sprints de duas semanas. É ele que ordena o product backlog e aceita ou recusa o que fica pronto.}
{Transformar pedidos soltos em histórias claras, com critérios de aceitação, ordenar o backlog com impacto e esforço à vista e montar o sprint backlog sem folhas de cálculo.}
{Pedidos mal descritos que obrigam a perguntar tudo de novo, estimativas que ninguém regista e perder o fio entre o que se pediu e o que a equipa entregou.}
{Utilizador principal da priorização, do backlog, das user stories e dos critérios de aceitação.}

\persona{Sofia Ferreira}{29 anos, responsável de operações numa empresa cliente}
{Usa o produto todos os dias e vai notando o que falta. Hoje manda os pedidos por email ao seu contacto na empresa fornecedora e raramente sabe o que lhes aconteceu.}
{Submeter um pedido em poucos minutos, ver se alguém já pediu o mesmo, votar nos pedidos que a afetam e saber em que estado está o que pediu.}
{Não ter resposta, não saber se o pedido foi lido e ter de repetir o mesmo pedido de cada vez que muda o interlocutor.}
{Cliente que submete, vota e acompanha. É também a nossa aproximação ao \textit{utilizador que vota}.}

\persona{Diogo Martins}{26 anos, programador full-stack}
{Está há dois anos na equipa. Recebe tarefas pelo Jira e por mensagens. Gosta de trabalhar com requisitos claros e de saber o motivo do que está a construir.}
{Ver a funcionalidade que lhe atribuíram com a história, os critérios de aceitação e as dependências, comentar com quem pediu, e ligar o que fez ao pedido original.}
{Pedidos vagos, mudanças de âmbito sem aviso e não saber que outra funcionalidade tem de estar pronta antes da sua.}
{Utilizador da atribuição de funcionalidades, do estado, das menções e da rastreabilidade.}

\persona{Helena Costa}{45 anos, gestora de projeto}
{Acompanha duas equipas e os prazos de cada release. Quer perceber a carga de cada equipa e se o plano se mantém, sem ter de interromper as pessoas para perguntar.}
{Ver o progresso do roadmap, a carga das equipas e o tempo que as funcionalidades levam do pedido à entrega, para ajustar prazos e avisar a direção a tempo.}
{Relatórios feitos à mão que já nascem desatualizados e descobrir os atrasos tarde demais.}
{Utilizadora principal dos dashboards, das métricas e dos milestones.}

\section{Cenários}

Cada cenário é o caminho de uma persona a tentar um objetivo. Inclui o que o sistema responde e o que acontece quando corre mal.

\subsection{C1. Sofia submete e acompanha um pedido}
A Sofia repara que os relatórios do produto não se exportam para Excel. Entra no sistema, escreve o título e a descrição do pedido e escolhe a categoria \textit{Relatórios}. Enquanto escreve, o sistema mostra dois pedidos parecidos. Um deles é o que ela queria: vota nele em vez de submeter outro. Dias depois recebe uma notificação a dizer que o pedido passou a \textit{Em análise}.

\textit{Quando corre mal.} Se o pedido que ela submete for mesmo duplicado e ela ignorar o aviso, o pedido fica criado mas o Ricardo pode juntá-lo ao original. A Sofia é avisada e os votos passam para o pedido principal.

\subsection{C2. Ricardo prioriza o backlog}
No início da semana o Ricardo abre a lista de pedidos novos. Filtra os que têm mais votos, regista para cada um o impacto e o esforço e vê a sugestão de prioridade da IA. Concorda com duas sugestões e discorda de uma: mantém a ordem dele e escreve a razão. Passa os pedidos escolhidos para o product backlog.

\textit{Quando corre mal.} Se a IA não estiver disponível, o sistema avisa e o Ricardo prioriza à mão. A IA só sugere, nunca decide. Se a estimativa de esforço ainda não existir, o pedido não entra no sprint backlog.

\subsection{C3. Marta monta uma release no roadmap}
A Marta abre o roadmap visual, cria a release 2.0 com um milestone para o fim de novembro e arrasta para lá as funcionalidades aprovadas. Ao arrastar \textit{Exportar relatórios}, o sistema avisa que depende de \textit{Filtros avançados}, que ainda não está agendada. Ela agenda as duas.

\textit{Quando corre mal.} Se criar uma dependência circular (A depende de B e B depende de A), o sistema recusa e explica porquê.

\subsection{C4. Ricardo gera histórias a partir de um pedido}
O pedido \textit{Exportar relatórios} foi aprovado. O Ricardo pede ao sistema para gerar as user stories. Aparecem três propostas no formato \textit{como\ldots{} quero\ldots{} para\ldots}. Edita uma, rejeita outra e acrescenta critérios de aceitação às que ficam. O sistema liga as histórias ao pedido original.

\textit{Quando corre mal.} Se o pedido tiver pouca informação, o sistema gera menos histórias e diz o que falta. As histórias geradas nunca entram no backlog sem o Ricardo as aceitar.

\subsection{C5. Diogo trabalha numa funcionalidade}
O Diogo vê no seu espaço de equipa a funcionalidade que lhe atribuíram, com as histórias e os critérios. Muda o estado para \textit{Em desenvolvimento}, escreve um comentário com uma pergunta e menciona a Sofia para esclarecer um detalhe. No fim liga o seu \textit{pull request} à funcionalidade e passa-a a \textit{Em revisão}.

\textit{Quando corre mal.} Se tentar passar a \textit{Concluída} sem pull request ligado, o sistema avisa que falta a ligação à implementação.

\subsection{C6. Helena verifica o andamento}
Antes da reunião com a direção, a Helena abre o dashboard. Vê o progresso do roadmap, a carga de cada equipa e o lead time e o cycle time médios. Repara que uma equipa está acima da capacidade e que a release 2.0 está atrasada. Exporta a vista e pede ao Ricardo para rever o âmbito.

\textit{Quando corre mal.} Se faltarem dados (por exemplo funcionalidades sem estimativa), o dashboard mostra a métrica como incompleta e diz quantos itens faltam, em vez de mostrar um número que engana.

\section{User stories}

Formato das aulas: como \textit{persona}, quero \textit{objetivo}, para \textit{benefício}. Cada história indica a área do enunciado, a prioridade proposta (MoSCoW) e os critérios de aceitação. As histórias são do nosso backlog, para construir este sistema. A geração automática de histórias (US-18) é uma funcionalidade do produto, não se confunde com este documento.

\subsection{Feature Management}

\us{US-01}{Feature Management}{cliente}{submeter um pedido de funcionalidade com título, descrição e categoria}{o meu pedido ficar registado e não se perder num email}{Must}{
\item O pedido só se grava com título e descrição preenchidos.
\item Depois de gravar, o pedido aparece na lista com o estado \textit{Novo}.
\item Quem submete recebe a confirmação no ecrã.}

\us{US-02}{Feature Management}{product manager}{categorizar e etiquetar os pedidos}{os agrupar por tema e encontrá-los depois}{Must}{
\item É possível atribuir uma categoria e várias etiquetas a cada pedido.
\item Pode filtrar-se a lista por categoria e por etiqueta.}

\us{US-03}{Feature Management}{product manager}{editar um pedido e ver as versões anteriores}{perceber o que mudou e quem mudou}{Should}{
\item Cada edição cria uma versão nova com autor e data.
\item É possível ver a lista de versões e comparar duas.
\item Voltar a uma versão antiga cria uma versão nova, não apaga o histórico.}

\us{US-04}{Feature Management}{utilizador}{pesquisar e filtrar pedidos por texto, estado, categoria e etiqueta}{encontrar depressa o que procuro}{Must}{
\item A pesquisa por texto procura no título e na descrição.
\item Os filtros combinam-se entre si.
\item Quando não há resultados, o ecrã diz isso e sugere limpar os filtros.}

\us{US-05}{Feature Management}{cliente}{ser avisado de pedidos parecidos quando escrevo o meu}{votar num que já existe em vez de criar um duplicado}{Should}{
\item Ao escrever o título, o sistema sugere pedidos semelhantes.
\item O product owner pode juntar dois pedidos; os votos e os comentários passam para o principal.}

\subsection{Prioritization}

\us{US-06}{Prioritization}{Product Owner}{ordenar manualmente os pedidos por prioridade}{a equipa trabalhar primeiro no que tem mais valor}{Must}{
\item Só o Product Owner altera a ordem.
\item A alteração fica registada com data e autor.
\item O Product Owner pode escrever a razão da ordem.}

\us{US-07}{Prioritization}{utilizador}{votar nos pedidos que me interessam}{mostrar o que mais precisa de ser feito}{Must}{
\item Cada utilizador tem um voto por pedido e pode retirá-lo.
\item O número de votos é visível na lista e no pedido.
\item O sistema impede votar duas vezes no mesmo pedido.}

\us{US-08}{Prioritization}{Product Owner}{receber uma sugestão de prioridade assistida por IA}{apoiar a minha decisão com outra leitura dos dados}{Should}{
\item A sugestão usa dados do pedido (votos, impacto, esforço).
\item A sugestão é só uma recomendação: nada muda sem o Product Owner aceitar.
\item Quando a IA falha ou não responde, o resto da priorização continua a funcionar.}

\us{US-09}{Prioritization}{Product Owner}{estimar o impacto e o esforço de cada pedido}{comparar pedidos e escolher melhor}{Must}{
\item Impacto e esforço usam uma escala fixa e documentada.
\item A lista pode ordenar-se por impacto, por esforço ou pela relação entre os dois.}

\subsection{Product Backlog Management}

\us{US-10}{Product Backlog}{Product Owner}{criar e manter o product backlog a partir dos pedidos aprovados}{ter uma lista única do que falta fazer}{Must}{
\item Um pedido aprovado pode ser acrescentado ao product backlog.
\item O backlog tem uma ordem e o Product Owner pode alterá-la.}

\us{US-11}{Product Backlog}{Product Owner}{definir o sprint backlog escolhendo itens do product backlog}{a equipa saber o que fazer em cada sprint}{Must}{
\item Um item só entra no sprint se tiver estimativa.
\item O sprint tem datas e objetivo.
\item O que não fica pronto volta ao product backlog.}

\us{US-12}{Product Backlog}{programador}{atualizar o estado de uma funcionalidade}{todos verem em que ponto está}{Must}{
\item Os estados são \textit{Novo}, \textit{Em análise}, \textit{Aprovado}, \textit{Rejeitado}, \textit{Em desenvolvimento}, \textit{Em revisão} e \textit{Concluído}.
\item Cada mudança de estado fica registada com data e autor.
\item Os interessados são notificados quando o estado muda.}

\us{US-13}{Product Backlog}{project manager}{atribuir funcionalidades a equipas de desenvolvimento}{cada trabalho ter um responsável}{Must}{
\item Uma funcionalidade pode ter uma equipa atribuída.
\item A lista pode filtrar-se por equipa.}

\subsection{Product Roadmap}

\us{US-14}{Product Roadmap}{product manager}{ver o roadmap numa linha do tempo visual}{perceber de relance o que vem e quando}{Must}{
\item O roadmap mostra releases e funcionalidades ao longo do tempo.
\item Arrastar uma funcionalidade para outra release atualiza-a.}

\us{US-15}{Product Roadmap}{product manager}{planear releases com data e funcionalidades}{comunicar o que entra em cada versão}{Must}{
\item Uma release tem nome, data prevista e lista de funcionalidades.
\item Uma funcionalidade pertence, no máximo, a uma release.}

\us{US-16}{Product Roadmap}{product manager}{definir milestones}{marcar os marcos importantes do produto}{Should}{
\item Um milestone tem nome, data e funcionalidades associadas.
\item O estado do milestone calcula-se a partir das funcionalidades.}

\us{US-17}{Product Roadmap}{programador}{ver e registar dependências entre funcionalidades}{saber o que tem de estar pronto antes do meu trabalho}{Should}{
\item O sistema avisa quando se agenda uma funcionalidade antes daquela de que depende.
\item O sistema recusa dependências circulares.}

\subsection{Requirements Engineering}

\us{US-18}{Requirements Engineering}{Product Owner}{gerar user stories a partir de um pedido aprovado}{poupar tempo a escrever e não perder nada do que foi pedido}{Should}{
\item As histórias saem no formato \textit{como\ldots{} quero\ldots{} para\ldots}.
\item O Product Owner edita, aceita ou rejeita cada uma.
\item Cada história aceite fica ligada ao pedido de origem.}

\us{US-19}{Requirements Engineering}{project manager}{seguir a ligação entre pedido, user story e implementação}{provar de onde veio cada coisa feita}{Should}{
\item A partir de um pedido vê-se as histórias e a implementação ligadas (por exemplo o pull request).
\item A partir de uma implementação chega-se ao pedido original.}

\us{US-20}{Requirements Engineering}{Product Owner}{escrever e gerir critérios de aceitação em cada história}{a equipa saber quando uma história está pronta}{Must}{
\item Cada história pode ter vários critérios.
\item Cada critério marca-se como cumprido ou não cumprido.
\item Uma história só passa a \textit{Concluído} se todos os critérios estiverem cumpridos, ou se o Product Owner justificar a exceção.}

\subsection{Collaboration}

\us{US-21}{Collaboration}{utilizador}{comentar e discutir um pedido}{esclarecer dúvidas no sítio onde está o pedido}{Must}{
\item Os comentários mostram autor e data.
\item O autor pode editar ou apagar o seu comentário.}

\us{US-22}{Collaboration}{utilizador}{receber notificações do que me interessa}{não ter de andar a verificar}{Should}{
\item Recebo notificação quando alguém me menciona, responde ou muda o estado de um pedido que sigo.
\item Posso escolher quais notificações quero receber.}

\us{US-23}{Collaboration}{programador}{mencionar colegas com @nome nos comentários}{chamar a pessoa certa para a conversa}{Should}{
\item Ao escrever @, aparece a lista de pessoas do espaço de trabalho.
\item A pessoa mencionada recebe uma notificação.}

\us{US-24}{Collaboration}{project manager}{ter um espaço de trabalho por equipa}{cada equipa ver só o seu trabalho}{Should}{
\item Cada espaço de trabalho tem membros, pedidos e funcionalidades.
\item Um utilizador só vê os espaços a que pertence.}

\subsection{Dashboards and Analytics}

\us{US-25}{Dashboards}{product manager}{ver quantos pedidos foram submetidos, aprovados e rejeitados}{perceber o fluxo de pedidos}{Should}{
\item Os totais podem filtrar-se por período.
\item Os números batem com a lista de pedidos.}

\us{US-26}{Dashboards}{product manager}{ver quantas funcionalidades há por release}{avaliar o tamanho de cada versão}{Should}{
\item O gráfico mostra o número de funcionalidades de cada release.
\item Carregar numa release leva à lista das suas funcionalidades.}

\us{US-27}{Dashboards}{project manager}{ver o lead time e o cycle time}{perceber quanto tempo demora uma funcionalidade do pedido à entrega}{Should}{
\item Lead time: da data do pedido à conclusão.
\item Cycle time: do início do desenvolvimento à conclusão.
\item O dashboard mostra a média e a evolução por período.}

\us{US-28}{Dashboards}{project manager}{ver a carga de cada equipa}{distribuir melhor o trabalho}{Should}{
\item A carga mostra o esforço estimado atribuído a cada equipa.
\item Quando a carga passa a capacidade definida, aparece um aviso.}

\us{US-29}{Dashboards}{project manager}{ver o progresso do roadmap}{saber se as releases estão dentro do plano}{Should}{
\item O progresso de cada release calcula-se a partir do estado das suas funcionalidades.
\item Uma release em atraso fica destacada.}

\subsection{Acesso}

\us{US-30}{Acesso}{utilizador}{registar-me e entrar no sistema com o meu perfil}{ver só o que o meu papel permite}{Must}{
\item Existem os perfis cliente, product manager, Product Owner, programador e project manager.
\item Cada perfil só acede às funcionalidades que lhe dizem respeito.
\item A palavra-passe não se guarda em texto simples.}

\subsection{Resumo}

\begin{center}
\begin{tabular}{lccc}
\toprule
Área & Must & Should & Total \\
\midrule
Feature Management & 3 & 2 & 5 \\
Prioritization & 3 & 1 & 4 \\
Product Backlog & 4 & 0 & 4 \\
Product Roadmap & 2 & 2 & 4 \\
Requirements Engineering & 1 & 2 & 3 \\
Collaboration & 1 & 3 & 4 \\
Dashboards & 0 & 5 & 5 \\
Acesso & 1 & 0 & 1 \\
\midrule
Total & 15 & 15 & 30 \\
\bottomrule
\end{tabular}
\end{center}

Nenhuma história está como \textit{Could} ou \textit{Won't} para já. Essa decisão é do Product Owner.

\section{Requisitos não funcionais}

Valores propostos e por confirmar com a equipa. O enunciado pede arquitetura justificada, modularidade e escalabilidade.

\begin{itemize}
\item \textbf{Usabilidade.} Submeter um pedido não deve exigir mais do que um ecrã. A interface funciona em portátil e em telemóvel.
\item \textbf{Desempenho.} A pesquisa e as listas respondem em menos de 2 segundos com alguns milhares de pedidos.
\item \textbf{Segurança.} Autenticação, perfis com permissões por papel, palavras-passe com hash e protegidas, e proteção contra os erros mais comuns (injeção, XSS).
\item \textbf{Rastreabilidade e auditoria.} Mudanças de estado, de prioridade e de versão ficam registadas com autor e data.
\item \textbf{Privacidade.} Só se guardam os dados pessoais necessários e cumpre-se o RGPD.
\item \textbf{Fiabilidade.} Se a IA ou um serviço externo falhar, o resto do sistema continua a funcionar.
\item \textbf{Manutenção.} Código modular, testes automáticos e integração contínua na \texttt{main}, como o contrato diz.
\item \textbf{Escalabilidade.} Arquitetura que permita crescer em utilizadores e em equipas sem refazer tudo.
\item \textbf{Acessibilidade.} Contraste, navegação por teclado e textos alternativos.
\item \textbf{Idioma.} Interface em português; a estrutura deve permitir acrescentar inglês.
\end{itemize}

\section{Interfaces candidatas}

O enunciado pede pelo menos 8 interfaces para além do login, do registo e do CRUD básico. Estas são as candidatas. Os wireframes do Sprint 2 escolhem e desenham as principais.

\begin{enumerate}
\item Roadmap visual com releases e milestones.
\item Quadro de priorização, com votos, impacto, esforço e sugestão da IA.
\item Detalhe do pedido, com versões, comentários e menções.
\item Gerador e editor de user stories com critérios de aceitação.
\item Planeamento de sprint (product backlog e sprint backlog).
\item Vista de rastreabilidade (pedido, história, implementação).
\item Dashboard de métricas (pedidos, lead time, cycle time, release).
\item Carga das equipas.
\item Centro de notificações.
\item Espaço de trabalho de equipa.
\end{enumerate}

\section{De onde vem cada persona}

\begin{center}
\begin{tabular}{lp{9cm}}
\toprule
Persona & Histórias principais \\
\midrule
Marta (product manager) & US-02, US-14, US-15, US-16, US-25, US-26 \\
Ricardo (Product Owner) & US-06, US-08, US-09, US-10, US-11, US-18, US-20 \\
Sofia (cliente) & US-01, US-05, US-07, US-21, US-22 \\
Diogo (programador) & US-12, US-17, US-23 \\
Helena (project manager) & US-13, US-19, US-24, US-27, US-28, US-29 \\
Todos & US-03, US-04, US-30 \\
\bottomrule
\end{tabular}
\end{center}

\section{Apresentação dos Sprints 1 e 2}

A apresentação oral do Sprint 1 e a do Sprint 2 fazem-se juntas. É o mesmo produto e o mesmo tema. Não são dois trabalhos (combinado da equipa, não vem do enunciado).

Para o Sprint 1 já temos: produto, problema, papéis, sete áreas, personas, cenários, user stories com critérios, requisitos não funcionais e interfaces candidatas. Para o Sprint 2 ficam os wireframes e a escolha da arquitetura. Falta o levantamento com o Character.ai.

\section{Pendente e limites}

\begin{itemize}
\item \textbf{Character.ai.} Falta o levantamento. Quando houver, confirmar personas, cenários e histórias.
\item \textbf{Validação pela equipa.} Ninguém além de quem escreveu leu este documento.
\item \textbf{Prioridades.} São uma proposta. O Product Owner decide.
\item \textbf{Escalas e números.} A escala de impacto e esforço, os tempos de resposta e os estados são valores propostos.
\item \textbf{O enunciado.} A página 3 é um esquema do percurso da cadeira (SWEBOK v4 e SFIA 9), não uma lista de ecrãs. As datas do enunciado (21/09 a 02/10 e seguintes) não são as do grupo, que estão no contrato e no Jira.
\item \textbf{Lab03.} O Smart Campus (\texttt{ENGENHARIA DE SOFTWARE.txt}) é o exercício da aula, usado só para o método. Nenhuma persona nem história dali entra aqui.
\item \textbf{Terminologia.} O enunciado chama Team Captain ao Scrum Master. No contrato e na equipa continua Scrum Master.
\end{itemize}

\end{document}
